\documentclass[11pt]{article}
\usepackage[margin=.78in,top=.65in,bottom=.65in]{geometry}
\usepackage{amsmath,amssymb,amsthm,hyperref}
\renewcommand{\normalsize}{\fontsize{10}{12}\selectfont}
\DeclareMathSizes{10}{10}{8}{6}
\hypersetup{colorlinks=true,linkcolor=blue,urlcolor=blue,citecolor=blue}
\setlength{\emergencystretch}{2em}
\newtheorem{theorem}{Theorem}
\newtheorem{lemma}[theorem]{Lemma}
\newcommand{\R}{\mathbb R}
\newcommand{\A}{\mathcal A}
\begin{document}
\begin{center}
{\Large\bfseries Repeated integrals and smooth endpoint data at every order}\\[.5em]
Research proof and Lean formalization\\
October 27, 2026
\end{center}
\begin{abstract}
We prove the all-order endpoint-data conjecture of Burke, Haris and
Madhavendra. The Stieltjes measure of a continuous nonconstant increasing
highest derivative is atomless and nonzero. Its finitely many moments have
a strictly positive smooth representing density equal to one near both
endpoints. A direct separation argument proves this representation, and
repeated integration supplies the required endpoint derivatives. The
admissible endpoint-data set is an invertible linear image of an open
convex smooth-moment image. The full all-order source statement is
formally verified in Lean, including the concrete convexity and openness
proofs and every source, measure, kernel and endpoint bridge.
\end{abstract}

\section{The source conjecture}
For $n\ge0$, let $W_n\subset\R^{n+1}$ consist of the vectors
$b=(b_0,\ldots,b_n)$ for which some $f\in C^n[0,1]$ satisfies
\begin{equation}\label{W}
 f^{(j)}(0)=0,\qquad f^{(j)}(1)=b_j\quad(0\le j\le n),
 \qquad f^{(n)}\text{ is nondecreasing and nonconstant}.
\end{equation}
Endpoint derivatives are one-sided. Burke, Haris and
Madhavendra~\cite{BHM}, Conjecture~1.2, p.~2, ask whether the following
statement $(P_n)$ holds for every $n$. Their Theorem~A proves it for $n\le3$.

\begin{theorem}\label{main}
For every integer $n\ge0$, the set $W_n$ is open. Every $b\in W_n$ has
a witnessing $F\in C^\infty(\R)$ such that
\[
 F^{(n+1)}(x)>0\ (0<x<1),\qquad
 F^{(n+1)}(0)=F^{(n+1)}(1)=1,
 \qquad F^{(j)}(0)=F^{(j)}(1)=0\ (j>n+1).
\]
In fact $F^{(n+1)}$ can be chosen identically one on neighborhoods of
both endpoints.
\end{theorem}

The source's sequel~\cite{Burke}, Theorem~4.2, retains the corresponding
all-order endpoint statement as a hypothesis. We prove $(P_n)$ itself;
we do not claim a separate discovery for its conditional applications or
for the previously proved cases $n\le3$.

\section{The smooth moment image}
Put $p_n(t)=(1,t,\ldots,t^n)$, and write
$m(\mu)=\int p_n\,d\mu$ for the moment vector of a finite positive measure
supported on $[0,1]$. Let
\begin{equation}\label{A}
 \A_n=\left\{\int_0^1p_n(t)\rho(t)\,dt:
 \begin{array}{l}
 \rho\in C^\infty(\R),\ \rho>0\text{ on }[0,1],\\[-2pt]
 \rho=1\text{ near }0\text{ and near }1
 \end{array}\right\}.
\end{equation}
All neighborhoods in this definition are neighborhoods in $\R$.

\begin{lemma}\label{open}
$\A_n$ is nonempty, convex and open, and $0\in\overline{\A_n}$.
\end{lemma}
\begin{proof}
The density one proves nonemptiness. Convex combinations preserve positivity
and both endpoint germs. To prove openness at a vector represented by $\rho$,
fix a nonnegative nonzero smooth function $\chi$ compactly supported in
$(0,1)$. The matrix
\[
 G_{ij}=\int_0^1\chi(t)t^{i+j}\,dt\qquad(0\le i,j\le n)
\]
is positive definite: its quadratic form at a nonzero coefficient vector
$c$ is $\int\chi(t)(\sum_i c_it^i)^2dt>0$. Indeed $\chi$ is positive on
some open interval and a nonzero polynomial has only finitely many roots.
Hence $G$ is invertible. For $c$ in a sufficiently small neighborhood of
zero, the densities
\[
 \rho_c(t)=\rho(t)+\chi(t)\sum_{i=0}^n c_it^i
\]
remain positive on $[0,1]$ and retain the endpoint germs. This uses the
positive minimum of $\rho$ and a uniform bound on the finite perturbation.
Their moment vectors are $m(\rho\,dt)+Gc$, which contain a neighborhood.

For the closure assertion, take $0<\delta<1$ and smooth cutoffs
$0\le\chi_0,\chi_1\le1$, each equal to one near its endpoint, supported
in the radius-$\delta$ ball around that endpoint. Choose $\delta<1/2$
so each cutoff vanishes near the opposite endpoint. The density
\[
 \beta_\delta=\delta+(1-\delta)(\chi_0+\chi_1)
\]
belongs to the family in~\eqref{A}. Its mass on $[0,1]$ is at most
$\delta+4\delta=5\delta$. Every coordinate of its moment vector lies
between zero and $5\delta$, proving convergence to zero.
\end{proof}

\begin{lemma}\label{representation}
Every finite nonzero atomless positive measure $\mu$ supported on $[0,1]$
has $m(\mu)\in\A_n$.
\end{lemma}
\begin{proof}
Suppose $y=m(\mu)\notin\A_n$. Since $\A_n$ is open and convex,
separation provides a continuous linear functional $L$ such that
\begin{equation}\label{sep}
 L(z)<L(y)\qquad(z\in\A_n).
\end{equation}
The closure assertion gives $0\le L(y)$. The functional is nonzero, since
otherwise~\eqref{sep} contradicts nonemptiness.

We claim $L(p_n(t))\le0$ for every $t\in[0,1]$. If it is positive at
an interior point, continuity supplies a nonnegative smooth bump $\psi$
supported inside $(0,1)$ with
\[
 \int_0^1L(p_n(t))\psi(t)\,dt>0.
\]
Every $1+c\psi$, $c\ge0$, is an allowed density in~\eqref{A}. Its
$L$-moment tends to $+\infty$ as $c\to\infty$, contradicting~\eqref{sep}.
Continuity extends the claim to both endpoints.

The polynomial $q(t)=-L(p_n(t))$ is nonnegative on $[0,1]$ and nonzero:
its coefficients are the negatives of $L$ on the coordinate basis. Its
zero set is finite, hence $\mu$-null. Since $\mu$ has positive mass,
\[
 0<\int q\,d\mu=-L(m(\mu))=-L(y),
\]
contradicting $0\le L(y)$. This proves the lemma, including $n=0$.
\end{proof}

\section{Identification of the endpoint data}
Let $f$ witness $b\in W_n$ and set $g=f^{(n)}$. Then $g$ is continuous,
nondecreasing and nonconstant, and $g(0)=0$. Extend it constantly outside
$[0,1]$. Its Stieltjes measure $\mu=dg$ is finite, atomless, supported
on $[0,1]$, and has positive mass $g(1)$. For every $0\le j\le n$,
\begin{equation}\label{kernel}
 b_j=\frac1{(n-j)!}\int_0^1(1-t)^{n-j}\,d\mu(t).
\end{equation}
For $j=n$, this is $g(1)=\mu([0,1])$. For $k=n-j\ge1$, Taylor's integral
remainder with zero initial jets gives
\[
 b_j=\int_0^1\frac{(1-x)^{k-1}}{(k-1)!}g(x)\,dx.
\]
Substitute $g(x)=\mu([0,x])$ and interchange the integrals. The inner
integral from $t$ to one is $(1-t)^k/k!$, proving~\eqref{kernel}. No
derivative of the original $g$ is required.

Define the linear map $T_n:\R^{n+1}\to\R^{n+1}$ by
\begin{equation}\label{T}
 (T_ny)_j=\sum_{i=0}^n
 \frac{\binom{n-j}{i}(-1)^i}{(n-j)!}\,y_i,
 \qquad 0\le j\le n,
\end{equation}
with $\binom{k}{i}=0$ for $i>k$. Before reversing the output coordinates,
its matrix is lower triangular with diagonal $(-1)^k/k!$. Thus $T_n$
is invertible, and the binomial theorem gives $b=T_nm(\mu)$.
By Lemma~\ref{representation}, choose a density $\rho$ as in~\eqref{A}
with the same moments. Set
\begin{equation}\label{primitive}
 F(x)=\frac1{n!}\int_0^x(x-t)^n\rho(t)\,dt\qquad(x\in\R).
\end{equation}
Repeated differentiation gives $F^{(n+1)}=\rho$ and zero initial jets
through $n$. Equations~\eqref{kernel}--\eqref{T} give the terminal jets
$b_j$. The endpoint germs of $\rho$ give all higher endpoint conditions
in Theorem~\ref{main}.

Conversely, construction~\eqref{primitive} for every density in~\eqref{A}
produces an element of $W_n$: its $n$th derivative is strictly increasing.
Therefore
\[
 W_n=T_n(\A_n).
\]
Lemma~\ref{open} and invertibility of $T_n$ prove openness, completing
Theorem~\ref{main}.

\section{Prior work and formal status}
Smooth representation of interior truncated moment functionals is established
in di Dio~\cite{diDio}, Theorem~2.5 and Corollary~2.7 of arXiv v2; the
latter is Corollary~3.1 in the published paper. These results belong to an
earlier moment-theoretic framework, not new machinery here. Our direct proof
uses standard separation, polynomial root finiteness and smooth bumps.
The source-specific content is the atomless Stieltjes identification,
the endpoint-one density requirement, and the all-order endpoint assembly.
The prior-work comparison is bounded. No absolute priority claim is made.

The current pinned Lean 4.34.1/mathlib \texttt{d13f23b7} development has
129 declarations in 18 modules. A fresh isolated replay treats warnings as
errors and checks every declaration's type and axioms; only
\texttt{propext}, \texttt{Classical.choice} and \texttt{Quot.sound} occur.
It formalizes the literal source hypotheses, Stieltjes measure and kernel
identity, invertible map~\eqref{T}, all derivatives of the smooth witness,
the endpoint baseline and the separation proof of
Lemma~\ref{representation}. A positive-definite bump Gram matrix proves
openness of the concrete moment image, and positivity-preserving density
combinations prove its convexity. The final endpoint
\begin{center}
\texttt{RepeatedIntegrals.source\_conjecture\_371768}
\end{center}
is closed: it quantifies over every natural order with no residual
convexity, openness, moment-representation or smoothness premises.


\begin{thebibliography}{9}
\bibitem{BHM} M. R. Burke, M. Haris and Madhavendra,
\emph{Repeated integrals of increasing functions},
\href{https://arxiv.org/abs/2512.02151v1}{arXiv:2512.02151v1} (2025).
\bibitem{Burke} M. R. Burke,
\emph{Comonotone approximation and interpolation by entire functions II},
\href{https://arxiv.org/abs/2512.23949v1}{arXiv:2512.23949v1} (2025).
\bibitem{diDio} P. J. di Dio,
\emph{Absolutely continuous representing measures of moment sequences
with an emphasis on the regularity of the densities},
\href{https://arxiv.org/abs/2411.03145v2}{arXiv:2411.03145v2} (2024).
Published truncated-moment version:
\href{https://doi.org/10.1007/s00013-025-02175-1}{doi:10.1007/s00013-025-02175-1}
(2025).
\end{thebibliography}
\end{document}
